\documentclass[11pt]{article}
\usepackage[T1]{fontenc}
\usepackage{textcomp}
\usepackage[margin=0.75in]{geometry}
\usepackage{amsmath,amssymb,amsthm}
\usepackage{array,booktabs}
\usepackage{hyperref}
\setlength{\parindent}{0pt}
\newtheorem{theorem}{Theorem}
\theoremstyle{definition}
\newtheorem{definition}{Definition}
\title{Flow Proof Artifact: prop-02-parallel-to-side-cuts-proportionally}
\author{Generated by \texttt{flow doc proof}}
\date{Flow Proof Book}
\begin{document}
\maketitle
If a straight line is drawn parallel to one side of a triangle, it cuts the other sides proportionally.\\[0.5em]
\textbf{Source.} euclid — Elements, Book VI, Proposition 2\\[0.5em]
\bigskip
\noindent{\large \textbf{Derived fact 1.} \textit{If a straight line is drawn parallel to one side of a triangle, it cuts the other sides proportionally} \hfill the Euclidean plane \cdot Euclid Book VI \cdot Proposition 2: a line parallel to a triangle side cuts the other sides proportionally}\par
\smallskip\noindent\textit{Coordinate: the Euclidean plane \cdot Euclid Book VI \cdot Proposition 2: a line parallel to a triangle side cuts the other sides proportionally}\par
\smallskip\noindent\textit{Source: euclid — Elements, Book VI, Proposition 2}\par
\smallskip\noindent\textit{Built on: proposition 1: triangles and parallelograms of the same height are as their bases, for Euclid Book VI on the Euclidean plane, proposition 37: triangles on the same base and between the same parallels are equal, for Book I of the Elements in on the Euclidean plane}\par
\medskip
\noindent\textbf{Goal.} If a straight line is drawn parallel to one side of a triangle, it cuts the other sides proportionally\par
\noindent$\displaystyle \text{parallel implies proportional sides}$\par
\medskip
\noindent
\begin{tabular}{@{} >{\raggedright\arraybackslash}p{0.06\textwidth} >{\raggedright\arraybackslash}p{0.47\textwidth} >{\raggedleft\arraybackslash}p{0.41\textwidth} @{}}
\textbf{\#} & \textbf{Proof} & \textbf{Mathematics} \\
\midrule
\textbf{1.} & We prove that proposition 2: a line parallel to a triangle side cuts the other sides proportionally for Euclid Book VI the Euclidean plane. & \\[0.45em]
\textbf{2.} & Let ABC = triangle. & \\[0.45em]
\textbf{3.} & Let DE = line\_parallel\_to\_BC. & \\[0.45em]
\textbf{4.} & Let AD = segment\_on\_AB. & \\[0.45em]
\textbf{5.} & From step 2, step 3, and step 4, we can deduce that ratio(AD, DB) equals ratio(AE, EC). & $\displaystyle ratio(AD, DB) = ratio(AE, EC)$ \\[0.45em]
\textbf{6.} & We invoke the derived fact governing Euclid Book VI the Euclidean plane: proposition 1: triangles and parallelograms of the same height are as their bases, for Euclid Book VI on the Euclidean plane. & \\[0.45em]
\textbf{7.} & We invoke the derived fact governing Book I of the Elements in the Euclidean plane: proposition 37: triangles on the same base and between the same parallels are equal, for Book I of the Elements in on the Euclidean plane. & \\[0.45em]
\textbf{8.} & From step 6 and step 7, this implies parallel implies proportional sides. Hence proven. & $\displaystyle \text{parallel implies proportional sides}$ \\[0.45em]
\bottomrule
\end{tabular}
\label{thm:Geometry-Euclid Book VI-proposition_2_a_line_parallel_to_a_triangle_side_cuts_the_other_sides_proportionally}
\par\medskip\hrule
\end{document}
